\section{\textbf{Experience}}

\cventry{September 2023 -- Present}{Data Engineer}{Division of Computational Pathology, \\ Department of Pathology and Laboratory Medicine, \\Indiana University School of Medicine \& IU Health}{Indianapolis, IN, U.S.A}{}{
\begin{itemize}
    \item Large Language Model 
    \item Engineered a comprehensive End-to-End pipeline for Whole Slide Images encompassing Tagging, Ingestion, Conversion, Anonymization, and Quality Control with Role-Based Access Control (RBAC) systems.
    \item Conducted extensive research and analysis for the selection of optimal Pathology Scanners, significantly enhancing the department's diagnostic capabilities.
    \item Spearheaded the development and deployment of advanced Visualization tools, optimizing data interpretation and analysis processes.
    \item Innovated a Low-Resource Optimization pipeline for efficient and accurate clinical inference, tailored for resource-limited clinic environments.
\end{itemize}
}

\cventry{September 2019 -- September 2023}{Data Analyst}{CBICA, Perelman School of Medicine,\\ University of Pennsylvania}{Philadelphia, PA}{U.S.A}{Under the supervision of Prof. Dr. S. Bakas
\begin{itemize}
    \item Pioneered a user-friendly Deep Learning Framework for Medical Imaging, integrating cutting-edge Digital Pathology features, requiring no programming expertise for end users.
    \item Architected an end-to-end pipeline for Digital Pathology, streamlining Ingestion, Anonymization, Preprocessing, and Segmentation/Classification of Whole Slide Images.
    \item Implemented a Deep Learning algorithm for Brain Extraction, vital to a global Tumor Segmentation federation, validated over 10,000 cases, enhancing pre-processing efficiency.
    \item Devised innovative strategies for optimizing Deep Learning algorithms in low-resource settings, notably enhancing medical imaging capabilities in underserved communities.
    \item Led the deployment of a clinical Deep Learning pipeline for Multiple Sclerosis, significantly reducing Gadolinium usage and cutting runtime by 90\%, thereby demonstrating substantial clinical improvements.
    \item Developed and validated a Brain Extraction software using Deep Learning, backed by comprehensive statistical analysis, ensuring high robustness and accuracy.
\end{itemize}
}

\cventry{January 2019 -- August 2019}{Research Intern}{CBICA, Perelman School of Medicine, University of Pennsylvania}{Philadelphia, PA}{U.S.A}{Under the supervision of Prof. Dr. C. Davatzikos and Prof. Dr. S. Bakas
\begin{itemize}
    \item Conducted in-depth analysis and evaluation of Brain Extraction techniques in the presence of diffuse gliomas.
    \item Developed and innovated a Deep Learning-based method for Brain Extraction, resulting in a deployable package, BrainMaGe.
\end{itemize}
}

\cventry{August 2018 -- September 2018}{Research Intern}{COESIP, Department of Electronics and Telecommunication, \\ SGGSIE\&T}{Nanded, MH}{India}{Under the supervision of Prof. Dr. S.N. Talbar and Dr. Ujjwal Baid
\begin{itemize}
    \item Secured 3rd position among 390 global research labs in the Brain Tumor Segmentation Challenge (BRaTS) 2018 at MICCAI, excelling in Survival Prediction of brain tumors.
    \item Implemented a Deep Learning-based workflow involving co-registration and Bias Field Correction using 3D-UNet for Brain Tumor Segmentation (BraTS).
    \item Engineered a post-segmentation pipeline for radiomic feature extraction in gliomas, and developed Machine Learning algorithms for Brain Tumor Survival Prediction.
\end{itemize}
}

\cventry{April 2018 -- July 2018}{Research Intern}{COESIP, Department of Electronics and Telecommunication,\\ SGGSIE\&T}{Nanded, MH}{India}{Under the supervision of Prof. Dr. S.N. Talbar
\begin{itemize}
    \item Developed Deep Learning and Computer Vision algorithms for the early detection of pulmonary lung nodules in lung cancer patients.
    \item Conducted extensive data preprocessing and analysis on a large dataset of lung CT scans, significantly enhancing the accuracy of the developed models.
    \item Created an innovative method for lung field extraction, substantially reducing the Region of Interest and boosting the model accuracy.
\end{itemize}
}

\cventry{May 2017 -- July 2017}{Research Intern}{HPC Lab, Department of Computer Science, SGGSIE\&T}{Nanded, MH}{India}{Under the supervision of Prof. H. P. Ambulgekar
\begin{itemize}
    \item Developed a Recommender System utilizing PySpark to implement Collaborative Filtering for movie recommendations.
\end{itemize}
}
